\documentclass[12pt,a4paper]{article}

\usepackage[margin=2cm]{geometry}
\usepackage[T2A]{fontenc}
\usepackage[utf8]{inputenc}
\usepackage[russian]{babel}
\usepackage{amsmath,amssymb,mathtools}
\usepackage{array,booktabs,longtable}
\usepackage{xcolor}
\usepackage{hyperref}
\usepackage{tikz}

\hypersetup{
  colorlinks=true,
  linkcolor=blue!55!black,
  urlcolor=blue!55!black,
  pdfauthor={Лёвин Артём Александрович},
  pdftitle={Ревью домашней работы}
}

\setlength{\parindent}{0pt}
\setlength{\parskip}{0.65em}
\setlength{\tabcolsep}{3pt}
\renewcommand{\arraystretch}{1.25}

\newcommand{\studentanswer}[1]{\textit{#1}}
\newcommand{\signaturefooter}{%
  \begin{tikzpicture}[remember picture,overlay]
    \node[anchor=south,font=\small] at ([yshift=7mm]current page.south)
    {Лёвин Артём Александрович эксклюзивно для Ксении Васильченко};
  \end{tikzpicture}%
}

\begin{document}

% -------------------- Титульный лист --------------------
\thispagestyle{empty}
\vspace*{\fill}
\begin{center}
  {\Huge\bfseries Ревью домашней работы}\par
  \vspace{1.4cm}
  {\large Автор: Лёвин Артём Александрович}\par
  \vspace{0.6cm}
  {\large 16 сентября 2026 года}\par
\end{center}
\vspace*{\fill}
\begin{center}
\begin{tikzpicture}[x=1cm,y=1cm]
  \draw[line width=0.7pt] (0,0) .. controls (2,0.7) and (4,-0.7) .. (6,0);
\end{tikzpicture}
\end{center}
\newpage

% -------------------- Сводная таблица --------------------
\section*{Сводная таблица}
\small
\begin{longtable}{>{\centering\arraybackslash}p{0.8cm} p{3.1cm} p{4.7cm} p{2.4cm} p{3.1cm}}
\toprule
\textbf{№} & \textbf{Тема} & \textbf{Суть ошибки} & \textbf{Ваш ответ} & \textbf{Правильный ответ} \\
\midrule
\endfirsthead
\toprule
\textbf{№} & \textbf{Тема} & \textbf{Суть ошибки} & \textbf{Ваш ответ} & \textbf{Правильный ответ} \\
\midrule
\endhead
\hyperref[task:3]{3} & Пересечение двух прямых & Решение системы не доведено до нахождения координат общей точки. & не указан & $(3;5)$ \\
\addlinespace
\hyperref[task:5]{5} & Вершина параболы & Дискриминант вычислен, но абсцисса вершины и её положение относительно оси $Oy$ не найдены. & не указан & $x_v=2$, справа от $Oy$ \\
\addlinespace
\hyperref[task:6]{6} & Пересечение прямой и параболы & Решение не начато: не составлено уравнение для общих точек. & не указан & $(0;1)$ и $(4;5)$ \\
\addlinespace
\hyperref[task:7]{7} & Параметр в уравнении прямой & При переносе единицы допущена ошибка со знаком, поэтому получено неверное значение параметра. & $m=-3$ & $m=3$ \\
\addlinespace
\hyperref[task:8]{8} & Горизонтальная прямая и парабола & Решение не начато; не использовано условие ровно одной общей точки. & не указан & $m=-1$ \\
\bottomrule
\end{longtable}
\normalsize

% -------------------- Разделитель --------------------
\newpage
\thispagestyle{empty}
\vspace*{\fill}
\begin{center}
  {\Huge\bfseries Разбор ошибок}
\end{center}
\vspace*{\fill}
\newpage

% -------------------- Задание 3 --------------------
\phantomsection\label{task:3}
\section*{Задание 3}

\textbf{Текст задания.}
Найдите точку пересечения прямых
\[
y=3x-4, \qquad y=-x+8.
\]

\textbf{Ваше решение.}
В тетради записаны исходные уравнения прямых, после чего решение не доведено до нахождения координат точки пересечения.

\textbf{Ваш ответ.}
\studentanswer{Не указан.}

\textcolor{red}{\textbf{Ошибка:}} решение системы не завершено.

\textbf{Где именно возникает ошибка.}
Последний достоверно правильный шаг --- запись двух уравнений прямых. Первый неполный шаг --- после этой записи не выполнено условие общей точки: не приравнены два выражения для одной и той же координаты $y$.

\textbf{Почему это неверно.}
Точка пересечения принадлежит обеим прямым одновременно. Поэтому в этой точке одно и то же значение $x$ должно давать одно и то же значение $y$ по обеим формулам. Пока два выражения для $y$ не приравнены, координата $x$ общей точки не найдена. Следовательно, нельзя получить и координату $y$, а значит, задача остаётся незавершённой.

\textcolor{blue!60!black}{\textbf{Ключевая идея:}}
для точки пересечения двух графиков значения функций при одном и том же $x$ равны. Поэтому нужно решить уравнение, полученное из равенства правых частей.

\textbf{Исправление от последнего правильного шага.}
Исходные уравнения уже записаны верно. Продолжаем именно с них:
\[
3x-4=-x+8.
\]
Переносим слагаемое $-x$ в левую часть, а число $-4$ --- в правую:
\[
3x+x=8+4.
\]
Получаем
\[
4x=12,
\]
откуда
\[
x=3.
\]
Теперь подставляем найденное значение $x$ в любое из исходных уравнений, например в первое:
\[
y=3\cdot 3-4=9-4=5.
\]

\textbf{Полное правильное решение.}
В точке пересечения обе прямые имеют одинаковые координаты, поэтому
\[
3x-4=-x+8.
\]
Решаем линейное уравнение:
\[
3x+x=8+4,
\]
\[
4x=12,
\]
\[
x=3.
\]
Находим ординату:
\[
y=3\cdot 3-4=5.
\]
Следовательно, общая точка имеет координаты
\[
(3;5).
\]

\textbf{Проверка.}
Проверим найденную точку по обеим формулам. Для первой прямой:
\[
3\cdot 3-4=5.
\]
Для второй прямой:
\[
-3+8=5.
\]
Обе формулы дают одинаковое значение $y=5$, значит, точка $(3;5)$ действительно лежит на обеих прямых.

\textcolor{teal!60!black}{\textbf{Итог:}}
точка пересечения прямых --- $(3;5)$.

\textbf{Задача на закрепление.}
Найдите точку пересечения прямых
\[
y=2x+1, \qquad y=-x+7.
\]

% -------------------- Задание 5 --------------------
\newpage
\phantomsection\label{task:5}
\section*{Задание 5}

\textbf{Текст задания.}
Для параболы
\[
y=2x^2-8x+3
\]
найдите абсциссу вершины и укажите, справа или слева от оси $Oy$ она находится.

\textbf{Ваше решение.}
Записано вычисление дискриминанта:
\[
D=64-24=40.
\]
Далее решение не продолжено.

\textbf{Ваш ответ.}
\studentanswer{Не указан.}

\textcolor{red}{\textbf{Ошибка:}} найденная величина не доведена до требуемого результата.

\textbf{Где именно возникает ошибка.}
Последний правильный шаг --- вычисление дискриминанта $D=40$. Само вычисление верно. Первый неполный шаг --- после него не найдена абсцисса вершины и не указано её положение относительно оси $Oy$.

\textbf{Почему это неверно.}
Дискриминант помогает исследовать корни квадратного уравнения и точки пересечения параболы с осью $Ox$. В этой задаче требуется другая характеристика --- абсцисса вершины. Для параболы
\[
y=ax^2+bx+c
\]
она находится по формуле
\[
x_v=-\frac{b}{2a}.
\]
Поэтому даже правильно вычисленный дискриминант сам по себе не отвечает на поставленный вопрос. Решение должно быть продолжено до нахождения $x_v$ и сравнения этого числа с нулём.

\textcolor{blue!60!black}{\textbf{Ключевая идея:}}
положение вершины по горизонтали определяется знаком её абсциссы: если $x_v>0$, вершина справа от оси $Oy$; если $x_v<0$, слева; если $x_v=0$, вершина лежит на оси $Oy$.

\textbf{Исправление от последнего правильного шага.}
Для данной параболы
\[
a=2, \qquad b=-8.
\]
По формуле абсциссы вершины:
\[
x_v=-\frac{-8}{2\cdot 2}=\frac{8}{4}=2.
\]
Так как
\[
2>0,
\]
вершина расположена справа от оси $Oy$.

\textbf{Полное правильное решение.}
Сравниваем уравнение
\[
y=2x^2-8x+3
\]
с общим видом
\[
y=ax^2+bx+c.
\]
Получаем
\[
a=2, \qquad b=-8, \qquad c=3.
\]
Абсцисса вершины равна
\[
x_v=-\frac{b}{2a}=-\frac{-8}{2\cdot 2}=2.
\]
Поскольку $x_v=2$ положительно, вершина находится справа от оси $Oy$.

\textbf{Проверка.}
Преобразуем выражение выделением полного квадрата:
\[
y=2(x^2-4x)+3.
\]
Добавим и вычтем $4$ внутри скобок:
\[
y=2\bigl((x-2)^2-4\bigr)+3.
\]
Тогда
\[
y=2(x-2)^2-8+3=2(x-2)^2-5.
\]
Из этой формы непосредственно видно, что вершина параболы имеет абсциссу $2$. Это подтверждает результат.

\textcolor{teal!60!black}{\textbf{Итог:}}
абсцисса вершины равна $2$; вершина расположена справа от оси $Oy$.

\textbf{Задача на закрепление.}
Для параболы
\[
y=3x^2+12x-5
\]
найдите абсциссу вершины и укажите, справа или слева от оси $Oy$ она находится.

% -------------------- Задание 6 --------------------
\newpage
\phantomsection\label{task:6}
\section*{Задание 6}

\textbf{Текст задания.}
Найдите точки пересечения прямой и параболы:
\[
y=x+1, \qquad y=x^2-3x+1.
\]

\textbf{Ваше решение.}
Решение не записано; возле номера задания поставлен знак вопроса.

\textbf{Ваш ответ.}
\studentanswer{Не указан.}

\textcolor{red}{\textbf{Ошибка:}} решение не начато.

\textbf{Где именно возникает ошибка.}
Последнего вычислительного шага, который можно продолжить, нет. Первый необходимый, но пропущенный шаг --- приравнять два выражения для $y$, потому что в общей точке прямая и парабола имеют одинаковые координаты.

\textbf{Почему это неверно.}
Чтобы найти точки пересечения двух графиков, недостаточно рассматривать каждое уравнение отдельно. Координаты общей точки должны удовлетворять обоим уравнениям одновременно. Поэтому нужно составить уравнение
\[
x+1=x^2-3x+1.
\]
Его корни дадут абсциссы всех общих точек. После этого для каждого найденного значения $x$ нужно вычислить соответствующее значение $y$.

\textcolor{blue!60!black}{\textbf{Ключевая идея:}}
точки пересечения графиков находятся из условия равенства значений функций при одном и том же $x$.

\textbf{Исправление.}
Приравниваем правые части:
\[
x+1=x^2-3x+1.
\]
Переносим всё в одну сторону:
\[
0=x^2-3x+1-x-1,
\]
то есть
\[
x^2-4x=0.
\]
Выносим $x$ за скобку:
\[
x(x-4)=0.
\]
Отсюда
\[
x=0 \quad \text{или} \quad x=4.
\]
Для $x=0$ по уравнению прямой
\[
y=0+1=1.
\]
Получаем точку $(0;1)$.

Для $x=4$:
\[
y=4+1=5.
\]
Получаем точку $(4;5)$.

\textbf{Полное правильное решение.}
В общей точке значения $y$ по обеим формулам равны, поэтому
\[
x+1=x^2-3x+1.
\]
Упрощаем:
\[
x^2-4x=0,
\]
\[
x(x-4)=0.
\]
Произведение равно нулю, если хотя бы один множитель равен нулю:
\[
x=0 \quad \text{или} \quad x=4.
\]
Находим ординаты по более простому уравнению $y=x+1$:
\[
x=0 \Rightarrow y=1,
\]
\[
x=4 \Rightarrow y=5.
\]
Следовательно, точки пересечения:
\[
(0;1) \quad \text{и} \quad (4;5).
\]

\textbf{Проверка.}
Для точки $(0;1)$:
\[
0+1=1,
\]
\[
0^2-3\cdot 0+1=1.
\]
Оба уравнения выполняются.

Для точки $(4;5)$:
\[
4+1=5,
\]
\[
4^2-3\cdot 4+1=16-12+1=5.
\]
Оба уравнения также выполняются.

\textcolor{teal!60!black}{\textbf{Итог:}}
прямая и парабола пересекаются в точках $(0;1)$ и $(4;5)$.

\textbf{Задача на закрепление.}
Найдите точки пересечения прямой и параболы:
\[
y=x+2, \qquad y=x^2-2x+2.
\]

% -------------------- Задание 7 --------------------
\newpage
\phantomsection\label{task:7}
\section*{Задание 7}

\textbf{Текст задания.}
При каком значении параметра $m$ точка $(2;7)$ лежит на прямой
\[
y=mx+1?
\]

\textbf{Ваше решение.}
Записано:
\[
7=2m+1,
\]
затем
\[
2m=-7+1,
\]
после чего получено
\[
m=-3.
\]

\textbf{Ваш ответ.}
\[
\studentanswer{m=-3.}
\]

\textcolor{red}{\textbf{Ошибка:}} неверное изменение знаков при переносе слагаемого.

\textbf{Где именно возникает ошибка.}
Последний правильный шаг:
\[
7=2m+1.
\]
Первый неверный шаг:
\[
2m=-7+1.
\]
Именно в этом переходе нарушается равносильность уравнения.

\textbf{Почему этот шаг неверен.}
Из равенства
\[
7=2m+1
\]
нужно убрать единицу из правой части. Для этого из обеих частей уравнения вычитают $1$:
\[
7-1=2m+1-1.
\]
Правая часть превращается в $2m$, а левая --- в $6$. Поэтому должно получиться
\[
6=2m,
\]
или, что то же самое,
\[
2m=6.
\]
В записанном переходе число $7$ необоснованно получило знак минус, а единица осталась со знаком плюс. Такое изменение не соответствует выполнению одной и той же операции над обеими частями уравнения.

\textcolor{blue!60!black}{\textbf{Ключевая идея:}}
при переносе слагаемого фактически выполняется одна и та же операция с обеими частями уравнения. Здесь нужно вычесть $1$ из обеих частей, а не менять знак у числа $7$.

\textbf{Исправление от последнего правильного шага.}
Начинаем с верной записи:
\[
7=2m+1.
\]
Вычитаем $1$ из обеих частей:
\[
7-1=2m,
\]
\[
6=2m.
\]
Делим обе части на $2$:
\[
m=3.
\]

\textbf{Полное правильное решение.}
Если точка $(2;7)$ лежит на прямой $y=mx+1$, её координаты должны удовлетворять уравнению прямой. Подставляем $x=2$ и $y=7$:
\[
7=2m+1.
\]
Вычитаем $1$ из обеих частей:
\[
6=2m.
\]
Делим на $2$:
\[
m=3.
\]

\textbf{Проверка.}
Подставляем $m=3$ и $x=2$ в уравнение прямой:
\[
y=3\cdot 2+1=6+1=7.
\]
Получено ровно то значение $y$, которое имеет заданная точка. Значит, параметр найден верно.

\textcolor{teal!60!black}{\textbf{Итог:}}
$m=3$.

\textbf{Задача на закрепление.}
При каком значении параметра $m$ точка $(4;13)$ лежит на прямой
\[
y=mx+1?
\]

% -------------------- Задание 8 --------------------
\newpage
\phantomsection\label{task:8}
\section*{Задание 8}

\textbf{Текст задания.}
Для какого значения $m$ горизонтальная прямая
\[
y=m
\]
имеет с параболой
\[
y=(x-2)^2-1
\]
ровно одну общую точку? Укажите это значение $m$.

\textbf{Ваше решение.}
Решение не записано; возле номера задания поставлен знак вопроса.

\textbf{Ваш ответ.}
\studentanswer{Не указан.}

\textcolor{red}{\textbf{Ошибка:}} не использовано условие ровно одной общей точки.

\textbf{Где именно возникает ошибка.}
Последнего вычислительного шага нет, потому что решение не начато. Первый необходимый шаг --- приравнять значения $y$ для прямой и параболы:
\[
m=(x-2)^2-1.
\]
Затем нужно использовать именно условие о количестве общих точек.

\textbf{Почему это неверно.}
Общие точки горизонтальной прямой и параболы соответствуют решениям уравнения
\[
m=(x-2)^2-1.
\]
После переноса числа $-1$ получаем
\[
(x-2)^2=m+1.
\]
Квадрат действительного числа неотрицателен. Если $m+1>0$, уравнение имеет два разных решения по $x$, а значит, две общие точки. Если $m+1<0$, действительных решений нет. Ровно одно решение возникает только тогда, когда правая часть равна нулю, потому что уравнение
\[
(x-2)^2=0
\]
имеет единственное решение $x=2$.

\textcolor{blue!60!black}{\textbf{Ключевая идея:}}
горизонтальная прямая имеет с параболой ровно одну общую точку тогда, когда проходит через вершину параболы. В данной записи вершина видна непосредственно: $(2;-1)$.

\textbf{Исправление.}
Приравниваем правые части:
\[
m=(x-2)^2-1.
\]
Переносим $-1$ в левую часть:
\[
m+1=(x-2)^2.
\]
Для единственного решения нужно
\[
m+1=0.
\]
Отсюда
\[
m=-1.
\]

\textbf{Полное правильное решение.}
Парабола записана в виде
\[
y=(x-2)^2-1.
\]
Так как квадрат $(x-2)^2$ принимает наименьшее значение $0$ при $x=2$, наименьшее значение функции равно
\[
y=-1.
\]
Следовательно, вершина параболы --- точка $(2;-1)$.

Горизонтальная прямая $y=m$ пересечёт параболу ровно один раз только тогда, когда пройдёт через её вершину. Поэтому
\[
m=-1.
\]

То же самое можно увидеть алгебраически:
\[
m=(x-2)^2-1,
\]
\[
(x-2)^2=m+1.
\]
Ровно один корень по $x$ получается при $m+1=0$, то есть при $m=-1$.

\textbf{Проверка.}
При $m=-1$ прямая имеет уравнение
\[
y=-1.
\]
Ищем общие точки:
\[
-1=(x-2)^2-1.
\]
Отсюда
\[
(x-2)^2=0,
\]
поэтому
\[
x=2.
\]
Получается ровно одна общая точка $(2;-1)$, что полностью соответствует условию.

\textcolor{teal!60!black}{\textbf{Итог:}}
$m=-1$.

\textbf{Задача на закрепление.}
Для какого значения $m$ горизонтальная прямая
\[
y=m
\]
имеет с параболой
\[
y=(x+1)^2+2
\]
ровно одну общую точку?

\signaturefooter
\end{document}
